\providecommand{\mergeind}[1]{\mathbf{1}_{\{#1\}}}
\section{Operand-free composition of whole-format indices}\label{merge:section}
\subsection{Record format and access contract}


Fix a binary format of precision $p\geq2$ with $E\geq1$ normal binades
and gradual underflow. Every addition is correctly rounded separately to
nearest, with ties to even. Inputs and operands are finite representable
numbers. The two signed zeros are identified numerically; NaNs, exception
flags, and zero-sign recovery are excluded. The outputs $-\infty$ and
$+\infty$ are absorbing under further finite additions. Scale values by
the subnormal spacing, so that this spacing is one. There are
\[
 N=2^p(E+1)-1
\]
finite numerical inputs, ordered by rank $0,\ldots,N-1$.

The whole-format index of Section~\ref{sec:full-format} partitions these ranks into
contiguous intervals refining $2E+1$ fixed dyadic input progressions.
On a finite-output piece of length $n$, the local inputs are $x_0+di$,
$0\leq i<n$, and the stored output is
\begin{equation}\label{merge:eq:record}
 f(i)=c+2H\left\lfloor\frac{i}{P}\right\rfloor
       +H\bigl(\mergeind{i\bmod P\geq t_1}+\mergeind{i\bmod P\geq t_2}\bigr),
 \qquad P=\frac{2H}{d},
\end{equation}
where $d,H$ are dyadic powers,
\[
 P\in\mathbb Z_{>0},\qquad 1\leq t_1\leq t_2\leq\min(P,n).
\]
The height $H=d/2$ is allowed. A threshold can be clipped to $n<P$;
the formula need only describe the stored interval. Infinite-output
pieces are constant. The full represented function is nondecreasing.
For the whole-format merge, a valid index covers every finite rank
without gaps, has finite outputs in the format, and uses heights from
$\{1/2,1,2,\ldots,2^{E-1}\}$. These invariants hold for the earlier
constructor and are preserved by the merge. The local closure lemma
itself does not need a bound on the height exponents.

Figure~\ref{merge:fig:record} illustrates a record produced by one
addition. In the format $p=4$, $E=2$, take the integer inputs $x=8+i$,
$0\leq i<8$, and the operand $8$. The input spacing is one, while the
outputs lie in the spacing-two binade $[16,32)$. In
\eqref{merge:eq:record}, this gives $c=16$, $H=2$, $P=4$, $t_1=2$,
and $t_2=3$. The unequal threshold positions encode the tie rule:
$17$ rounds to $16$, whereas $19$ rounds to $20$.

\begin{figure}[htbp]
\centering
\setlength{\unitlength}{1pt}
\begin{picture}(260,132)
\put(25,18){\vector(1,0){215}}
\put(25,18){\vector(0,1){99}}
\put(241,14){\makebox(12,10)[l]{$i$}}
\put(0,120){\makebox(235,10)[l]{$\RN(16+i)-16$}}
\multiput(25,16)(27,0){8}{\line(0,1){4}}
\put(22,4){\makebox(6,9){$0$}}
\put(49,4){\makebox(6,9){$1$}}
\put(76,4){\makebox(6,9){$2$}}
\put(103,4){\makebox(6,9){$3$}}
\put(130,4){\makebox(6,9){$4$}}
\put(157,4){\makebox(6,9){$5$}}
\put(184,4){\makebox(6,9){$6$}}
\put(211,4){\makebox(6,9){$7$}}
\put(10,34){\makebox(10,9){$2$}}
\put(10,54){\makebox(10,9){$4$}}
\put(10,74){\makebox(10,9){$6$}}
\put(10,94){\makebox(10,9){$8$}}
\multiput(23,38)(0,20){4}{\line(1,0){4}}
\put(25,18){\circle*{4}}
\put(52,18){\circle*{4}}
\put(79,38){\circle*{4}}
\put(106,58){\circle*{4}}
\put(133,58){\circle*{4}}
\put(160,58){\circle*{4}}
\put(187,78){\circle*{4}}
\put(214,98){\circle*{4}}
\end{picture}
\caption{An attained two-threshold record. Only integer query indices are
plotted. Translating the input index by $P=4$ translates the output by
$2H=4$. On the first period the two upward crossings occur at indices
$2$ and $3$.}
\label{merge:fig:record}
\end{figure}

The merge receives only these records and the public format. It does not
receive the operands or their rounding history. We count abstract exact
arithmetic and comparison operations. Integer bit costs are separate.

Finite-state transfer composition and tree-based prefix computation are
established tools \cite{ref1,ref2}. Likewise, pulling an interval
boundary back through a monotone map produces one cut; analogous
breakpoint propagation is classical for monotone piecewise-linear
composition \cite[Lemma~2.1]{ref14}. We use these facts
without claiming them as new. The question addressed here is whether the
particular dyadic records remain closed under composition, and whether
the resulting records can be constructed directly from two indices.
The generalized-inverse composition principle is also classical:
for nondecreasing periodic integer maps, Connes and Consani
\cite[Proposition~4.1]{connesconsani2013} explicitly give the inverse
threshold relation and its reverse-order composition law. The formulas
below specialize this principle to the stored rounding records.

\subsection{Explicit inversion and local closure}

\begin{lemma}[Inverse of one record]\label{merge:lem:inverse}
For a finite record \eqref{merge:eq:record}, let $\operatorname{inv}_f(z)$ be the first
$i\in\{0,\ldots,n-1\}$ with $f(i)\geq z$, or $n$ if none exists.
If $z\leq c$ then $\operatorname{inv}_f(z)=0$. Otherwise put
\[
 k=\left\lceil\frac{z-c}{H}\right\rceil,\qquad
 q=\left\lfloor\frac{k-1}{2}\right\rfloor,\qquad s=k-2q\in\{1,2\}.
\]
Then
\begin{equation}\label{merge:eq:inverse}
 \operatorname{inv}_f(z)=\min\{n,qP+t_s\}.
\end{equation}
\end{lemma}

\begin{proof}
In each period, the first crossing of level $c+(2q+1)H$ is $qP+t_1$,
and the first crossing of $c+(2q+2)H$ is $qP+t_2$. A threshold equal
to $P$ is the first index of the next period, giving the same crossing.
Coincident thresholds skip the middle level and give both crossings
at the same index. Every output lies in $c+H\mathbb Z$, so crossing $z$ is
equivalent to crossing $c+kH$.

If a threshold was clipped to $n<P$, extend the displayed formula using
that stored threshold. It remains monotone and agrees with the record
on $0\leq i<n$. Clipping its inverse at $n$ removes any artificial
crossing outside the interval. The case $P=1$, where both thresholds
equal one, is included. The difference $z-c$ also covers negative
anchors and targets.
\end{proof}

\begin{corollary}[Inverse queries]\label{merge:cor:inverse-query}
On a whole-format index with $q$ pieces, the first input rank whose
output is at least a finite real target $z$, or strictly greater than $z$,
can be found in $O(\log(q+1))$ abstract operations. The exact preimage
of a finite closed output interval $[a,b]$, $a\leq b$, is a half-open
input-rank interval computable in the same bound.
\end{corollary}

\begin{proof}
The last outputs of successive pieces are nondecreasing. Binary search
finds the first piece whose last output meets the requested inequality;
the explicit inverse then finds its first satisfying input. For a strict
inequality, return the piece start if $z<c$, and otherwise replace $k$
in Lemma~\ref{merge:lem:inverse} by $\lfloor(z-c)/H\rfloor+1$.
Constant infinite pieces are handled by their order relative to the
finite target. The preimage of $[a,b]$ runs from the first output at
least $a$ up to, but excluding, the first output strictly greater than
$b$. An absent crossing has rank $N$. In subnormal-spacing units each
finite output $y$ is an integer. For a doubled target $t=2z$, the two
tests are therefore $2y\geq\lceil t\rceil$ and
$2y\geq\lfloor t\rfloor+1$, respectively. These formulas apply to
nonintegral $t$ as well.
\end{proof}

\noindent\emph{Binary32 example.}
In ordinary numerical units let $f(x)=\RN(x+1)$ and take $z=1$.
The binary32 neighbours of $1$ are $1-2^{-24}$ and $1+2^{-23}$.
Both midpoints round to $1$ by the even tie rule, so $f(x)=1$
holds exactly for the binary32 inputs in $[-2^{-25},2^{-24}]$.
In ascending order of finite numerical inputs, counting from zero and
identifying the two zeros, the first rank with $f(x)\geq1$ is
$1\,283\,457\,023$. The first rank with $f(x)>1$ is $3\,003\,121\,664$.
The half-open rank interval thus contains $1\,719\,664\,641$ inputs.

\begin{lemma}[Local composition]\label{merge:lem:closure}
Let $A$ and $B$ be finite records of the form \eqref{merge:eq:record}.
Suppose the image under $A$ of a contiguous input interval $I$ lies
within the valid input interval of $B$. On $I$, the composite has the
same record form, with input spacing $d_A$ and height
\[
 H=\max(H_A,H_B).
\]
The assertion includes clipped thresholds, fictitious half-input-spacings,
and constant restrictions.
\end{lemma}

\begin{proof}
Extend the formula for $A$ to every integer local index. It is monotone,
has values in $c_A+H_A\mathbb Z$, and commutes with displacement $2H_A$ in
numerical input units.

For $B$, with numerical input origin $y_B$, define an artificial extension
to every real $y$. Set
\[
 u=\frac{y-y_B}{d_B},\qquad q=\left\lfloor\frac{u}{P_B}\right\rfloor,
 \qquad r=u-qP_B,
\]
and put
\[
 b_{\mathrm{ext}}(y)=c_B+2H_Bq+
 H_B\bigl(\mergeind{r\geq t_{1,B}}+\mergeind{r\geq t_{2,B}}\bigr).
\]
This agrees with the valid $B$ record, is nondecreasing, has range in
$c_B+H_B\mathbb Z$, and satisfies
\[
 b_{\mathrm{ext}}(y+2H_B)=b_{\mathrm{ext}}(y)+2H_B.
\]
At a period boundary its left limit is at most the next period's initial
value, including when thresholds coincide or equal $P_B$. Clipped
thresholds merely select a possibly artificial continuation.

Let $G=b_{\mathrm{ext}}\circ A$ on integer indices. Since $2H$ is an
integer multiple of both numerical periods, and $2H/d_A$ is a positive
integer,
\[
 G(i+2H/d_A)=G(i)+2H.
\]
If $H_B\geq H_A$, its values lie in $G(0)+H\mathbb Z$ by the range property
of $B$. If $H_A>H_B$, the dyadic relation implies $H_A$ is an integer
multiple of $2H_B$. All differences between $A$ values are multiples of
$H_A$, and the equivariance of $b_{\mathrm{ext}}$ preserves each such
difference exactly. The same coset property follows.

On one period, monotonicity and equivariance bound $G$ between $G(0)$
and $G(0)+2H$. The coset leaves only these two levels and $G(0)+H$.
Their first crossing indices are the two thresholds. Rebase at the first
index of $I$ and clip to its length. On $I$ both extensions agree with
the actual records, proving the assertion.
\end{proof}

\begin{remark}
The proof does not assume that the extension of $A$ lies everywhere on
the input lattice of $B$. Such an assumption would require special care
when a stored height is a fictitious half-input-spacing. The real
extension in the proof removes that requirement. No values outside a
valid record need to be evaluated by the algorithm.
\end{remark}

\subsection{A linear-work merge}

\begin{theorem}\label{merge:thm:merge}
Let $A,B$ be valid whole-format indices with $q_A,q_B$ pieces. Extend
both to the two infinities by absorption. From their records alone, one
can construct an index for $B\circ A$ with
\[
 q_{\mathrm{out}}\leq\min\{N,q_A+q_B-1\}
\]
pieces using $O(q_A+q_B)$ abstract arithmetic and comparison operations
and $O(q_A+q_B)$ auxiliary storage. The output has the same record
format, with fixed-width description size
\[
 O\bigl(q_{\mathrm{out}}(p+\log(E+1))\bigr).
\]
\end{theorem}

\begin{proof}
For every internal $B$-piece boundary, let $y_j$ be the numerical input
of the piece immediately to its right. Find the first finite input rank
$i$ with $A(i)\geq y_j$. Monotonicity gives one cut, possibly at $0$ or
$N$. Merge these at most $q_B-1$ cuts with the $q_A-1$ internal cuts
of $A$ and the two endpoints. Repeated cuts give no empty pieces. Thus
the number of nonempty intervals is at most $q_A+q_B-1$, and also
at most $N$.

Each resulting interval belongs to one $A$ record, and its finite image
lies in one $B$ record. An image jump can skip several $B$ pieces; the
corresponding preimage cuts then coincide. Infinite outputs of $A$ are
already separated by its own endpoints and remain infinite. If the
applicable $B$ piece is infinite, the composite is constant there.
Otherwise Lemma~\ref{merge:lem:closure} applies.

For construction, sweep the sorted boundary values $y_j$ through the
last outputs of the sorted $A$ pieces. Advance past a piece when its
last output is below the current boundary. Use Lemma~\ref{merge:lem:inverse}
on the first remaining piece. Each pointer advances only linearly many
times, so all preimage cuts cost $O(q_A+q_B)$ operations. Merge the
two sorted cut lists in linear work. The numerical input origins of
$B$ pieces are computed by one sweep through the fixed input
progressions. These number $2E+1\leq q_B$, because their boundaries
are retained.

Now sweep the new intervals and their parent records. For a finite
interval $[l,r)$, evaluate $c=B(A(l))$, put
\[
 H=\max(H_A,H_B),\qquad P=2H/d_A,\qquad T=\min(P,r-l),
\]
and recover the two thresholds. For target $z=c+jH$, first invert the
parent $B$ record. If its inverse returns the stop sentinel, the new
threshold is $T$. Otherwise convert that input rank to its numerical
value $y$ using the parent's input origin and spacing, invert the
parent $A$ record at $y$, and clip the resulting rank to $[l,l+T]$.
Subtracting $l$ gives the new local threshold. This is exact: on the
interval, $B(A(i))\geq z$ holds precisely when $A(i)$ has reached the
first $B$ input whose output is at least $z$. Two applications of the
explicit inverse suffice for each target. Thus each new interval costs
constant work after its parents have been identified.

For the bit bound, rank endpoints and finite anchors use
$O(p+\log(E+1))$ bits. The dyadic spacings and heights have $O(E)$
possible exponents. Taking a maximum introduces no new exponent range.
Thresholds are clipped to a piece length smaller than $2^p$, because
the initial progression boundaries are retained. Constant infinities
require only a tag. These fields give the asserted fixed-width bound.
\end{proof}

\begin{remark}
Composition is associative for the represented functions and admits
the identity. The theorem does not require a unique encoding or
byte-for-byte equality under different merge parenthesizations.
\end{remark}

\subsection{Balanced construction}

\begin{lemma}\label{merge:lem:leaf}
A one-addition whole-format index can be built in $O(E)$ abstract
operations with at most $\min(N,4E+1)$ pieces.
\end{lemma}

\begin{proof}
Begin with the $2E+1$ initial input progressions. Pull back the
$2(E-1)$ interior rounding-region boundaries and the two overflow
thresholds under $x\mapsto x+a$, and merge the sorted cuts linearly.
For the negative overflow threshold $-\tau$, use the first exact sum
strictly greater than $-\tau$; for positive overflow use the first
exact sum at least $\tau$. Thus both overflow ties belong to the
infinite tails. These cuts add at most $2E$ pieces.

On each finite interval the operation is one uniform rounder
$\RN_h(x+a)$. For a requested output level $z$, put $k=\lceil z/h\rceil$.
The lower boundary of the exact inputs whose rounded output is at least
$z$ is $h(k-1/2)$. This boundary is included if $k$ is even and excluded
if $k$ is odd.
Solving the resulting non-strict or strict inequality on the local
progression $x_l+di+a$ gives each threshold in constant operations.
The anchor is one rounded addition. Both the input and rounding-region
lists can be traversed once, giving total $O(E)$ work.
\end{proof}

\begin{corollary}\label{merge:cor:balanced}
For a block of $m\geq1$ finite operands, a whole-format index can be
constructed in
\[
 O\bigl(E(1+m\log(m+1))\bigr)
\]
abstract operations, with at most $\min(N,4Em+1)$ pieces, query cost
$O(\log(E(m+1)+2))$, and description length
\[
 O\bigl(\min(N,4Em+1)(p+\log(E+1))\bigr).
\]
For $m=0$, the $2E+1$-piece identity costs $O(E)$ to initialize.
\end{corollary}

\begin{proof}
Build the $m$ leaves with Lemma~\ref{merge:lem:leaf} and compose adjacent
blocks in a balanced binary tree. Induction using
Theorem~\ref{merge:thm:merge} gives at most $4Es+1$ pieces at a node
containing $s$ operands. Each level has total record count $O(Em)$,
and there are $O(\log(m+1))$ levels. This gives the work bound.
Discarding consumed levels gives $O(Em)$ storage. The query and bit
bounds follow from the record representation.
\end{proof}

The piece constant is larger than the earlier sequential constructor's
$2E+1+2Em$. The gain is an operand-free merge and an improved abstract
preprocessing bound, not a smaller guaranteed final index. With sequential
linear merges, the critical-path sum is $O(Em)$; the corollary therefore
does not establish sublinear parallel depth or a GPU speedup. In the
scaled-integer implementation, values can occupy $O(p+E)$ bits, so
the arithmetic-operation count is not a machine-word cost bound.

\subsection{An obstruction for a restricted query grid}

\begin{theorem}\label{merge:prop:obstruction}
For $p\geq6$, put $R=2^p$ and
\[
 X_\ast=\{R/2+8+4i:0\leq i<R/16-2\}.
\]
There is no general composition rule determined only by the restrictions
of two block transfer functions to $X_\ast$, even when their trajectories
satisfy the original two-binade corridor promise.
\end{theorem}

\begin{proof}
Take $A=(1)$, $B=(3/2,-3/2)$, and $C=(0,0)$. These operands are
representable under the original format assumptions. Both $B$ and $C$
have real sum zero and return zero from zero. Every $x\in X_\ast$ is
divisible by four, lies at most $3R/4-4$, and all relevant exact and
rounded states remain in the lower unit-spaced binade. Ties-to-even gives
\[
 \RN(x+3/2)=x+2,\qquad
 \RN(x+2-3/2)=\RN(x+1/2)=x.
\]
Thus $F_B|_{X_\ast}=F_C|_{X_\ast}$ is the identity. But $F_A(x)=x+1$,
and
\[
 \RN(x+1+3/2)=\RN(x+5/2)=x+2,
 \qquad F_B(F_A(x))=x,
\]
whereas $F_C(F_A(x))=x+1$. A rule receiving only the restrictions cannot
distinguish $B$ and $C$, yet their composites with $A$ differ at every
query. This is a contradiction.
\end{proof}

The proposition concerns summaries determined only by the restricted
function. It does not exclude encodings that preserve additional
off-grid behavior or the original operands. Whole-format indices retain
the required behavior and cover absorbing infinities, which explains
why Theorem~\ref{merge:thm:merge} has a different conclusion.
